
\begin{proposition}[The same non-uniruled components occur with every even Euler number]\label{cdp:st:C.2.UN.14}
Fix a proper holomorphic family $p^\times:\mathcal T^\times\to\Delta^\times$
of connected complex two-tori. A marked smooth filling is a proper
holomorphic map $f:Y\to\Delta$ from a smooth complex threefold with
connected fibres, together with a biholomorphism
$f^{-1}(\Delta^\times)\simeq\mathcal T^\times$ over $\Delta^\times$.
For each irreducible component $S_i$ of $(Y_0)_{\rm red}$, let
$Z_i\to S_i^\nu\to S_i$ be a smooth compact resolution of its
normalization. Write $\mathcal E(f)$ for the finite multiset of
bimeromorphism classes $[Z_i]$ for which $Z_i$ is not uniruled.
Here a smooth compact surface is uniruled when it admits a dominant
meromorphic map $C\times\PP^1\dashrightarrow Z$ for a smooth compact
curve $C$. Repetitions in the multiset record distinct components.

For every integer $k$, the following equality holds between sets of
finite multisets:
\[
 \begin{split}
 \{\mathcal E(f): f\text{ is a marked smooth filling of }p^\times\}
 ={}&\{\mathcal E(f): f\text{ is such a filling},\\[-2pt]
 &\hspace{48mm}e((Y_0)_{\rm red})=2k\}.
 \end{split}                                                     \tag{C.EC.1}\label{cdp:eq:C.EC.1}
\]
For each filling in the left-hand set, the representative on the
right is obtained by the explicit proper holomorphic vertical
blowup map of \ref{cdp:st:C.2.UN.4}. That map induces a bijection between the
original non-uniruled components and the resulting non-uniruled
components, and a bimeromorphism of their normalization resolutions.
Every original full-fibre multiplicity is preserved by this bijection.

The same assertion holds for compact smooth total spaces over a
fixed compact base curve, with the marked family fixed away from a
finite set $Q$. One can prescribe an arbitrary even Euler number at
each point of $Q$ simultaneously, preserving the multiset and
multiplicity at each point. In particular, the collection of
non-uniruled surface classes which occur in arbitrary compact
torus-fibred fillings is exactly the collection which occur in
compact fillings all of whose reduced fibres have Euler number zero.
\end{proposition}
\begin{proof}
We specify the component comparison for one of the smooth blowups
$\pi:Y'\to Y$ used in \ref{cdp:st:C.2.UN.4}. Its centre $C$ is a point or a
smooth compact curve in the central support. The inverse image of
that support consists of the strict transforms $S_i'$ of its
original components together with the exceptional divisor
\[
 E=\PP(N_{C/Y})\longrightarrow C.                    \tag{C.EC.2}\label{cdp:eq:C.EC.2}
\]
Each restriction $S_i'\to S_i$ is a proper holomorphic map and is an
isomorphism over the complement of $C$. Since $C$ has smaller
dimension than $S_i$, this complement is dense. Thus this actual
map is bimeromorphic. Normalizing both surfaces and resolving them
gives a bimeromorphism $Z_i'\dashrightarrow Z_i$: explicitly, take
the closure of the graph of the isomorphism on the dense normal
open sets, normalize it, and resolve the resulting compact analytic
surface. Its two proper holomorphic projections resolve the stated
meromorphic map. The pullback of the defining equation of the full
fibre has its original vanishing order at the generic point of
$S_i'$, because $\pi$ is an isomorphism there. Hence its
multiplicity is unchanged.

Uniruledness is preserved by these bimeromorphisms. Indeed, composing
a dominant meromorphic map $C_0\times\PP^1\dashrightarrow Z_i$
with either bimeromorphic inverse gives such a map for $Z_i'$; its
dominance follows on the dense open set where that inverse is an
isomorphism. Resolving the graph makes the composition a
well-defined meromorphic map. The same argument in the reverse
direction proves the equivalence.

The newly introduced divisor in \eqref{cdp:eq:C.EC.2} is uniruled. For a point
centre it is $\PP^2$, which is rational: the map
$\PP^1\times\PP^1\dashrightarrow\PP^2$ given on the affine chart
by $(u,v)\mapsto[1:u:v]$ is bimeromorphic. For a curve centre it is
a projective line bundle over $C$. Its rank-two bundle is
meromorphically trivial on $C$: tensoring by a sufficiently positive
line bundle supplies two global sections independent at the
generic point, and their ratio gives the meromorphic
trivialization. Projectivizing it gives a bimeromorphism
$C\times\PP^1\dashrightarrow\PP(N_{C/Y})$.
Its subsequent strict transforms remain uniruled by the same
composition argument. Therefore a finite sequence of these blowups
introduces no new non-uniruled class and loses none of the original
classes. It preserves their repetitions and multiplicities as well.

Starting from any marked filling, \ref{cdp:st:C.2.UN.4} constructs a finite
sequence of exactly these maps whose final reduced Euler number is
$2k$. Every centre lies over zero, so the blowdown is a
biholomorphism over $\Delta^\times$ and transports the given marking
by composition. We have just proved that its multiset is the
original $\mathcal E(f)$. This gives inclusion of the left side of
\eqref{cdp:eq:C.EC.1} in the right. Inclusion in the other direction follows
because each filling with Euler number $2k$ is among the same
marked fillings defining the left side. This also proves the
assertion when there are no fillings: both sets are empty.

For the compact version, apply the same construction successively
at the finitely many points of $Q$. Every centre is a compact
analytic submanifold of the current compact threefold. Its global
blowup is compact and smooth, and the composite with the fixed base
map is proper holomorphic with connected fibres, by the
structure-sheaf and Stein-factor comparison in \ref{cdp:st:C.2.UN.4}. Its
restriction off that point is an isomorphism. Thus a construction
at one point preserves every fibre and every marking already
treated at another point. There are only finitely many critical
values of a compact proper map to a curve which is smooth
generically: the image of its compact analytic critical locus is
a proper analytic subset of the curve. Taking $Q$ to contain them
all and taking $k=0$ at each of them proves the final assertion;
the remaining smooth torus fibres already have Euler number zero.
\end{proof}

For completeness, a surface dominated meromorphically by
$C\times\PP^1$ has no nonzero pluricanonical section. The pullback
of such a section is nonzero, because a dominant meromorphic map
between smooth surfaces is generically of full rank. It is
holomorphic away from the finitely many indeterminacy points and
extends over them by Hartogs' theorem. Its restriction to any
$\{c\}\times\PP^1$ has degree $-2m$ in pluricanonical degree $m>0$
and must vanish. This contradicts its being nonzero. Thus the
equality applies in particular to a specified general-type
normalization class. It makes the occurrence problem for such a
component part of the unrestricted zero-Euler problem through the
explicit maps just proved. It does not supply a comparison between
two independently specified marked fillings. The construction
uses \ref{cdp:st:C.2.UN.4} only and invokes no assertion that such a comparison
extends across the central fibre.


